\section{The walk}
\label{sec:walk}

Pollard's rho~\cite{pollard1978} in the parallel collision-search
form of van Oorschot and Wiener~\cite{vanoorschot1999} walks many
pseudo-random trails through the group, reports the distinguished
points it passes, and solves the logarithm when two trails meet. On a
Koblitz curve the Frobenius $\sigma$ and negation generate a group of
automorphisms of order $2m=262$ that act on $\langle P\rangle$, and a
walk defined on the orbits of that group needs $\sqrt{262}\approx16.2$
times fewer steps~\cite{wienerzuccherato1998,glv2000,duursma1999}.
The expected number of classes visited before a collision on
$N=(\ell-1)/262$ classes is $\sqrt{\pi N/2}$, which is $2^{60.809}$
here.\art{ecc2k130/WALK-CONSTANT.md, §1} Every walk in this paper is
a map on those classes; what differs between them is how the class
action is respected and how much each step costs.

\subsection{The $\sigma$-walk}
\label{sec:walk-sigma}

The campaign walks the map of Bailey et al.~\cite{bailey2009}:
\begin{align*}
  j &= 3 + \bigl((\HW(x_R) \gg 1) \bmod 8\bigr) \in [3,10],\\
  R' &= \sigma^j(R) + R,\\
  \text{distinguished} &\iff \HW(x_R)\le w,
\end{align*}
with $w=34$ in the client's default and $w=32$ in the live campaign.
Because the weight is a class invariant, $f(\sigma R)=\sigma f(R)$ and
$f(-R)=-f(R)$, so the map descends to classes without any canonical
representative being computed inside the loop. The division by two
in the selector is load-bearing: $x$-coordinates of subgroup points
have even weight, and a selector on the raw weight would use only
four of the eight branches.\art{gpu/ecc2k/README.md}

One step is one affine addition: one inversion, two multiplications
and a free Frobenius, since $\sigma^j$ is wiring. As a scalar the step
is $R\mapsto[1+s^j]R$ where $s$ is the root of $s^2+s+2=0$ that
$\sigma$ realises on $\langle P\rangle$. The $\sigma^j$ selection is
branch-free: three conditional squarings, each one \texttt{LOP3} per
word, and $j$ is never materialised.\art{ecc2k130/README.md, ``Iteration function''}

No coefficients are tracked. A walk is identified by its seed, and a
report is the pair (seed, canonical endpoint), 32 bytes. The start
point is $R_0 = Q + \sum_i c_i\,\sigma^i(P)$ with $c$ a 128-bit string
from a PRF of the seed, computed with the same arithmetic and merged
into finished lanes under a mask; tracking a linear combination in
the loop would cost a 129-bit modular multiplication per step, which
this walk cannot afford. Collision resolution re-walks both trails
counting how often each branch $1+s^j$ was applied, which gives
$\text{endpoint} = [\mu](\alpha_0 P + Q)$ with $\mu=\prod_j(1+s^j)^{n_j}$,
matches the two endpoints up to Frobenius and negation, and solves
for $k$ modulo $\ell$.\art{ecc2k130/README.md, ``Distinguished points and restarts''}
The $\sigma$-walk has no fruitless cycles: $1+s^j$ is never a root of
unity, so a trail cannot return to its own point.

The client can also carry the eight branch counters as it goes
(\texttt{WITNESS=1}). That is eight counters and not a coefficient,
because the product $\mu$ commutes and nothing is multiplied per step.
It turns a report into something a third party checks with one
double scalar multiplication, about $200$ group operations, instead
of re-walking the $2^{25.27}$ steps a weight-34 report represents, an
asymmetry of $2^{17.4}$. On the bitsliced CPU walk it costs $6.0\%$
measured ($40.026$ to $37.762$~M~it/s, against $33.8\%$ predicted); on
the packed GPU walk it cost $34.7\%$ ($14.003$ to $9.150$~B/s), above
the $2\%$ the note had set as its abandonment threshold, so the
campaign builds pin \texttt{WITNESS=0} and the record stays 32
bytes.\art{ecc2k130/CAIRN-WITNESS.md}\art{ecc2k130/benchmarks/witness-cost/}
A later fused kernel narrows the packed cost ($9.347867$ to
$11.800903$~B/s in the witness-carrying mode) without closing
it.\art{ecc2k130/benchmarks/sigma-fused/RESULTS.md}

\subsection{The $r$-adding table walk}
\label{sec:walk-table}

The $\sigma$-walk pays for $\sigma^j$ in the basis conversions around
it on a GPU (Section~\ref{sec:arith-packed}). An $r$-adding walk on
the same classes removes that cost:
\begin{align*}
  h(R) &= (\HW(x)/2) \bmod 8,\\
  k(R) &= \Bigl(\textstyle\sum_e L(e)\,x_e\Bigr)\cdot \HW(x)^{-1} \bmod 131,\\
  \varepsilon(R) &= \text{bit $p(R)$ of } y,\\
  R' &= R + (-1)^{\varepsilon}\,\sigma^{k}(T_h),
\end{align*}
with $T_h=[a_h]P+[b_h]Q$ and a $131\times8$-point table ($48.7$~KB) in
shared memory. $k$ shifts by one under $\sigma$ and $\varepsilon$
flips under negation, so the walk descends to classes; the step is
one affine addition with a table operand whose Frobenius power is a
lookup rather than a conversion.\art{ecc2k130/ITERATION-FUNCTION.md, §4}

Unlike the $\sigma$-walk, an adding walk on classes has fruitless
cycles~\cite{bls2011}. Their source here is $\tau^2+\tau+2=0$: four-step
patterns in which the table points' Frobenius powers satisfy a
$\tau$-relation, and six-step pairwise patterns. The first cycle rule
let 24 four-step and 4 six-step patterns through; the count observed
against the count predicted was $1{,}145$ against $1{,}312$
($0.87\pm0.03$) for the $\tau$-relation class and 126 against 166 for
the pairwise class.\art{ecc2k130/WALK-CONSTANT.md, §5} At the
campaign's trail length of $2^{28.41}$ steps those cycles trapped
$53.7\%$ of trails, a loss of $8.634\times$ with the $2^{30}$ iteration
guard. A revised rule (v2) brings residual trapping to
$2.0\cdot10^{-5}$ of trails at a $1.0048\times$ iteration cost from
merge parting; a third revision (v3, Section~\ref{sec:walk-cycles})
replaces history-dependent exits with a raw-cycle anchor.

\subsection{The $x$-only doubling-and-bridge map}
\label{sec:walk-bridge}

On the power-capped RTX~PRO~4500 a third map pays. The common
transition is $R' = R + \sigma(R)$, which as a scalar is $[1+s]$, and
on this curve $(1+s)$ has norm 2 and the step can be computed from
$x$ alone. That scalar alone does not generate the group, so a sparse
bridge $R' = R + \sigma^3(R)$ is taken when $\HW(x) \bmod 72 = 14$;
the two scalar multipliers together generate the full group, which
the earlier $[2]/[3]$ map had failed to do (``the scalar subgroup
defect'').\art{ecc2k130/README.md, ``ECC2K-130 and ECC2K-95''}
A frozen study of $2{,}000$ trials on $\mathrm{GF}(2^{23})$, where
both the modulus-32 control and the bridge predicate select only
weight 14, gives a mean-work ratio of $1.006696$ for the bridge map
against the control, so the map's throughput gain
(Section~\ref{sec:g7}) is not paid back in extra iterations. The map
passes an independent $2{,}048$-point affine
checkpoint.\art{ecc2k130/benchmarks/g7-12b-xonly-doubling-bridge/README.md}

\subsection{The walk constant, measured}
\label{sec:walk-constant}

Throughput is half of a rate. The other half is the number of
iterations a walk needs, which for a non-random map differs from the
random-mapping $\sqrt{\pi N/2}$ by a constant $c$. The repository
measures $c$ as iterations divided by a random mapping's iterations
on the class set, on toy Koblitz curves of the same family, with the
device walk and a host emulation measured separately and an
$r$-adding control with 256 branches alongside.\art{ecc2k130/WALK-CONSTANT.md, §4}

\begin{table}[t]
\centering
\small
\caption{Walk constant $c$ on toy Koblitz curves, ECC2K-130 branch
  distribution ($H=8$). Floor $c=1$. The adding-walk control with 256
  branches measures $1.0002\pm0.0012$ at $n=23$. The injective-branch
  model $1/\sqrt{1-\sum p^2}$ predicts $1.070$ for the $\sigma$-walk
  at $n=131$.}
\label{tab:walkconst}
\begin{tabular}{@{}rlll@{}}
\toprule
$n$ & measured on & table walk & $\sigma$-walk \\
\midrule
23 & host emulation & $1.0037\pm0.0012$ & $1.0995\pm0.0013$ \\
23 & device & $1.0076\pm0.0037$ & $1.0948\pm0.0040$ \\
31 & device & $1.0035$ & $1.0964$ \\
37 & device ($W=16$) & $1.0036$ & $1.0949$ \\
41 & device & $1.0125\pm0.0059$ & $1.0910\pm0.0064$ \\
59 & device & $0.9949$ & $1.0817$ \\
\midrule
\multicolumn{2}{@{}l}{pooled / extrapolated to $n=131$} & $1.0035\pm0.0006$ & $[1.070,\,1.082]$ \\
\bottomrule
\end{tabular}
\end{table}

The $\sigma$-walk is a one-to-one map on each branch and pays the
injective-branch penalty $1/\sqrt{1-\sum_j p_j^2}$, which with the
weight distribution at $n=131$ ($\sum p_j^2=0.127$) is $1.070$; the
measurements in Table~\ref{tab:walkconst} drift from $1.0995$ at $n=23$
toward it, and the extrapolation to $n=131$ is $[1.070,1.082]$. The
table walk is not injective per branch and sits at $1.0035\pm0.0006$
pooled. The ratio $c_t/c_\sigma$ is $0.928$--$0.938$. Expected work
for the campaign's $\sigma$-walk is therefore
\[
  2^{60.809}\times[1.070,\,1.082] = 2^{60.91}\text{--}2^{60.92}
\]
iterations, with a further $1.185\times$ (that is, $2^{61.15}$--$2^{61.17}$)
under the original $2^{30}$ restart guard, since a guard at four trail
lengths discards $15.6\%$ of steps. Bailey et al.'s $2^{60.9}$
stands.\art{ecc2k130/WALK-CONSTANT.md, §6} The guard was raised to
$2^{32}$, twelve trail lengths at weight 32, which discards $0.007\%$
of steps ($1.00007\times$).\art{ecc2k130/benchmarks/max-iters/README.md}

On $\mathrm{GF}(2^{41})$ with 128 walks, cutoff 13 and 64 planted
logarithms, the $\sigma$-walk needed $164{,}200\pm9{,}900$ iterations
per solve and the table walk $157{,}800\pm7{,}500$, ratio
$0.96\pm0.07$, against a bare prediction of $102{,}600$ with about
$1.5\times$ harness overhead from the distinguished-point
tail.\art{ecc2k130/ITERATION-FUNCTION.md, §6.2}

\subsection{Cost per solve, not rate}
\label{sec:walk-cycles}

Putting rate and constant together is where the table walk's story
changes. Its throughput advantage of Section~\ref{sec:tablewalk} is
$1.09$--$1.17\times$; its walk constant is $0.93\times$; but its
as-built cycle rule trapped half of all campaign-length trails, so
the cost per solve of the as-built kernel was $4.85$--$6.26\times$ the
$\sigma$-walk's. With rule v2 the projection is $0.81$--$0.86\times$,
and the switch deadline (the fraction of the campaign after which
switching no longer pays) is $15$--$22\%$. That projection has not been
tested: the paired rate of the v2 kernel that would test it is
declared in the note and has not been run.\art{ecc2k130/WALK-CONSTANT.md, §11}
The campaign therefore still walks the $\sigma$-map, and the
$20$~B/s table-walk rows of Section~\ref{sec:tablewalk} are rates, not
costs per solve.

Rule v3 then replaced the history-dependent exits with a raw-cycle
probe of at most eight steps and a least-key anchor exit, checked by
$19{,}200$ refused $\tau$-cycles in a dedicated test, a $4{,}096$-point
host probe with no mismatch and planted logarithms $8/8$; its scatter
check against the predicted cycle counts gives $\chi^2=8.79$ on 15
degrees of freedom (the v1 rule: $23.30$).\art{ecc2k130/CYCLE-ESCAPE-V3.md}
V3 requires a new table corpus; the campaign's $\sigma$-walk corpus is
unaffected.
